% !TeX root = ../peskin_main.tex

\section{Interacting Fields and the S-Matrix    }
\label{sec:interacting}

% (4.18)の計算
相互作用描像における時間発展の演算子
\[
    U(t, t_0) = e^{iH_0 (t - t_0)} e^{-iH(t-t_0)}
\]
の時間微分を求めよ。

% 解
\[
    \begin{aligned}
        i \frac{\partial}{\partial t} U(t, t_0)
         & = i \frac{\partial}{\partial t} \left( e^{iH_0 (t - t_0)} e^{-iH(t-t_0)} \right)                    \\
         & = i \left( iH_0 e^{iH_0 (t - t_0)} e^{-iH(t-t_0)} + e^{iH_0 (t - t_0)} (-iH) e^{-iH(t-t_0)} \right) \\
         & = e^{iH_0 (t - t_0)} (H - H_0) e^{-iH(t-t_0)}                                                       \\
         & = e^{iH_0 (t - t_0)} H_{\mathrm{int}} e^{- iH_0 (t - t_0)} e^{iH_0 (t - t_0)} e^{-iH(t-t_0)}        \\
         & = H_{I}(t) U(t, t_0)
    \end{aligned}
\]
ただし、
\[
    H_{I}(t) = e^{iH_0 (t - t_0)} H_{\mathrm{int}} e^{- iH_0 (t - t_0)}  = \int d^3 x \frac{\lambda}{4!} \phi_{I}^4(x)
\]。

\dotfill
% (4.29)の計算
\[
    i \frac{\partial}{\partial t} U(t, t_0) = H_{I}(t) U(t, t_0), \quad U(t_0, t_0) = 1, H_{I}(t) = \int d^3 x \frac{\lambda}{4!} \phi_{I}^4(x)
\]
より、\( U(t, t_0) \)を\(\lambda\)のべき級数で展開せよ。

% 解
初期条件 \( U(t_0, t_0) = 1 \) より、
\[
    U(t, t_0) = 1 - i\int_{t_0}^{t} dt' H_{I}(t') U(t', t_0)
\]
右辺の\(U\)に右辺全体を代入すると
\[
    \begin{aligned}
         U(t, t_0) 
         & = 1 - i\int_{t_0}^{t} dt' H_{I}(t') \left( 1 - i\int_{t_0}^{t'} dt'' H_{I}(t'') U(t'', t_0) \right) \\
         & = 1 - i\int_{t_0}^{t} dt' H_{I}(t') + (-i)^2 \int_{t_0}^{t} dt' \int_{t_0}^{t'} dt'' H_{I}(t') H_{I}(t'') U(t'', t_0)
    \end{aligned}
\]
これを延々と繰り返すと、
\[
    \begin{aligned}
         U(t, t_0) 
         & = 1 - i\int_{t_0}^{t} dt' H_{I}(t') + (-i)^2 \int_{t_0}^{t} dt' \int_{t_0}^{t'} dt'' H_{I}(t') H_{I}(t'') + \cdots \\
         & = 1 + \sum_{n=1}^{\infty} (-i)^n \int_{t_0}^{t} dt_1 \int_{t_0}^{t_1} dt_2 \cdots \int_{t_0}^{t_{n-1}} dt_n H_{I}(t_1) H_{I}(t_2) \cdots H_{I}(t_n)
    \end{aligned}
\]
が得られる。\(H_I\)は\(\lambda\)を一つ含むため、これが\(\lambda\)のべき級数展開であることがわかる。

\dotfill
% (4.22)の計算

\[
    U(t, t_0) = = 1 + \sum_{n=1}^{\infty} (-i)^n \int_{t_0}^{t} dt_1 \int_{t_0}^{t_1} dt_2 \cdots \int_{t_0}^{t_{n-1}} dt_n H_{I}(t_1) H_{I}(t_2) \cdots H_{I}(t_n)
\]
を、時間順序積を用いて表せ。

% 解
重積分の初歩を思い出すと、二重積分
\[
    \int_{a}^{b} dx \int_{a}^{x} dy f(x) f(y)
\]
は、\(x\) の範囲が \([a, b]\)、\(y\)の範囲が\([a, b]\)の正方形を線分\(x=y\)で二等分したうちの、\(x>y\)の部分（下三角領域）での積分である。

一方、\(x<y\)の部分（上三角領域）での積分は、
\[
    \int_{a}^{b} dy \int_{a}^{y} dx f(x) f(y)
\]
である。これらを足し合わせると、四角形全体での積分となる。
ここで、上式の積分変数\(x, y\)を入れ替えると
\[
    \int_{a}^{b} dx \int_{a}^{x} dy f(y) f(x)
\]
となり、元の積分から\(f(x), f(y)\)の順序が入れ替わった積分となる。
このことから、元の積分は
\[
    \int_{a}^{b} dx \int_{a}^{x} dy f(x) f(y) = \frac{1}{2} \int_{a}^{b} dx \int_{a}^{b} dy T \left\{f(x) f(y)\right\}
\]
と表すことができる。ただし、\(T\)は時間順序積であり、\(x, y\)を時間とみなしている。

これを\(n\)重積分に拡張すると、
\[
    \int_{t_0}^{t} dt_1 \int_{t_0}^{t_1} dt_2 \cdots \int_{t_0}^{t_{n-1}} dt_n H_I(t_1) H_I(t_2) \cdots H_I(t_n) = \frac{1}{n!} \int_{t_0}^{t} dt_1 \int_{t_0}^{t} dt_2 \cdots \int_{t_0}^{t} dt_n T \left\{H_I(t_1) H_I(t_2) \cdots H_I(t_n)\right\}
\]
となる。
したがって、\(U(t, t_0)\)は
\[
    U(t, t_0) = 1 + \sum_{n=1}^{\infty} \frac{(-i)^n}{n!} \int_{t_0}^{t} dt_1 \int_{t_0}^{t} dt_2 \cdots \int_{t_0}^{t} dt_n T \left\{H_I(t_1) H_I(t_2) \cdots H_I(t_n)\right\}
\]
と表すことができる。
さらに、右辺を指数関数のテイラー展開と見ることで
\[
    U(t, t_0) = T\left[ \exp \left( -i \int_{t_0}^{t} dt' H_I(t') \right) \right]
\]
となる。

\dotfill

% 